% discussion.tex -- Cross-Enterprise Verification Paper (RU-V7)

\section{Discussion}

\subsection{Hypothesis Evaluation}

The hypothesis---that a hub-and-spoke model can verify cross-enterprise
interactions with $< 2\times$ gas overhead while preserving data
privacy---is \textbf{confirmed} for the primary scenario (2--50 enterprises
with linear interaction density). All three verification approaches achieve
$< 2\times$ overhead in this regime: Sequential at 1.41$\times$, Batched
Pairing at 0.64$\times$, and Hub Aggregation at 1.16$\times$.

The null hypothesis---that cross-enterprise verification requires $\geq
2\times$ gas overhead, or that privacy breaks down, or that proof generation
exceeds 120 seconds---is \textbf{rejected} on all three criteria. The
overhead is well below 2$\times$; privacy is preserved with exactly 1 bit
leakage (interaction existence, which is the theoretical minimum); and
proof generation is estimated at 448~ms (rapidsnark), two orders of
magnitude below the 120-second threshold.

\subsection{Approach Trade-offs}

The three verification approaches offer distinct trade-offs suited to
different deployment scenarios:

\textbf{Sequential verification} requires no coordination between
enterprises. Each enterprise submits its proof independently, and
cross-reference proofs are submitted as separate transactions. The
overhead (1.41--1.81$\times$) is acceptable for enterprise counts up to 50.
However, sequential verification fails the $< 2\times$ target when
interaction density is high ($M \gg N$), limiting its applicability to
sparse interaction graphs.

\textbf{Batched pairing verification} achieves the lowest gas cost
(0.64$\times$ for 2 enterprises) by sharing pairing computations across all
proofs in a single transaction. This requires a hub coordinator to collect
proofs before submission, introducing a synchronization point. The approach
remains below 1$\times$ overhead in all tested configurations, including
the dense interaction case.

\textbf{Hub aggregation} (SnarkPack/Nebra UPA-style) offers a middle ground:
lower gas than sequential at the cost of an aggregation step. The per-proof
gas decreases as $O(1/N)$ due to amortization of the fixed aggregation
cost. Hub aggregation is the recommended approach for large-scale
deployments ($N > 30$) where batched pairing transaction size may exceed
block gas limits.

\subsection{Comparison with Production Systems}

\textbf{Rayls (Parfin/Drex)}: Rayls uses a hub chain for cross-enterprise
coordination with the Enygma protocol combining ZK proofs and homomorphic
encryption. While Rayls achieves strong privacy within subnets, its hub
chain sees transaction metadata (source, destination, amount commitments).
Our protocol's hub coordinator sees \emph{only} Groth16 proofs and
Poseidon commitments---strictly less information than Rayls' hub chain.
The trade-off is that Rayls supports richer cross-enterprise operations
(e.g., atomic transfers with HE) while our protocol is limited to
\emph{verification} of interactions that have already been recorded in both
enterprises' state trees.

\textbf{Polygon AggLayer}: AggLayer aggregates ZK proofs from multiple
chains with public state, achieving \$0.002--0.003 per transaction at scale.
Our system operates in a fundamentally different trust model (private
enterprise state), but at comparable per-proof gas costs when using batched
verification ($\sim$122K gas/proof at $N=2$, decreasing with scale).

\textbf{zkSync Gateway}: Similar to AggLayer, Gateway aggregates proofs
for cost reduction. At \$0.0047/tx with $\sim$3,895 tx/batch, Gateway
operates at a much larger scale than our enterprise-focused system (2--50
enterprises, 1--49 interactions per batch).

\subsection{Limitations}

\begin{enumerate}
\item \textbf{Stage 1 results only}: This paper reports Stage 1
  (Implementation) results. Stochastic baseline (Stage 2), adversarial
  testing (Stage 3), and ablation (Stage 4) are future work.

\item \textbf{Analytical gas model}: Gas costs are computed analytically
  rather than measured on-chain. While our model is validated against
  on-chain measurements (RU-V3: 205,600 gas predicted vs. 205,600 gas
  measured for Groth16 verification), edge cases in EVM gas accounting
  (memory expansion, call depth) may introduce deviations.

\item \textbf{Proof generation is estimated}: The 448~ms rapidsnark
  estimate is derived from the constraint count and a per-constraint
  timing ratio from RU-V2. Actual circuit compilation and proving with
  Circom/snarkjs may reveal additional overhead from witness generation
  or file I/O.

\item \textbf{Single interaction type}: We evaluate a simple
  ``matching value'' interaction predicate. More complex predicates
  (range proofs, conditional logic, multi-party interactions) would
  increase the constraint count and proving time.

\item \textbf{Groth16 trusted setup}: The cross-reference circuit
  requires its own Groth16 trusted setup (separate from enterprise
  state transition circuits). For $N$ different enterprise circuit
  configurations plus 1 cross-reference circuit, $N + 1$ trusted setups
  are needed.

\item \textbf{Privacy is per-interaction, not per-relationship}: The
  protocol does not hide the \emph{pattern} of interactions over time.
  Repeated cross-reference submissions between the same two enterprises
  reveal the interaction frequency, which may be sensitive in competitive
  markets.
\end{enumerate}

\subsection{Reconciliation with Literature}

All experimental metrics are consistent with published benchmarks:

\begin{itemize}
\item Groth16 verification gas matches Nebra's formula
  $(181 + 6L)$K within 0.3\%.
\item Merkle proof generation and verification times match RU-V1
  within 10\%.
\item Constraint counts follow the established
  $1{,}038 \times (d+1)$ formula from RU-V2 exactly.
\item The 1-bit-per-interaction leakage is consistent with the
  theoretical minimum for ledger-visible cross-reference protocols
  (any on-chain record of a cross-enterprise event inherently reveals
  the event's existence).
\end{itemize}

No published system achieves lower leakage for cross-enterprise
verification on a public ledger, confirming that our protocol is at
the information-theoretic boundary.
